\section{Numerical Comparison with Measurement}
\label{sec:numerical_validation}


\CatchFileBetweenTags{\AlphaInvVal}{constants.tex}{AlphaInvVal}
\CatchFileBetweenTags{\AlphaInvExperimentalVal}{constants.tex}{AlphaInvCodataExperimentalVal}
\CatchFileBetweenTags{\AlphaInvSigma}{constants.tex}{AlphaInvCodataSigma}

\CatchFileBetweenTags{\FermiConstSigma}{constants.tex}{FermiConstSigma}
\CatchFileBetweenTags{\AlphaSSigma}{constants.tex}{AlphasSSigma}
\CatchFileBetweenTags{\AlphaSExperimentalVal}{constants.tex}{AlphasSExperimentalVal}
\CatchFileBetweenTags{\HiggsMassSigma}{constants.tex}{HiggsMassSigma}

\CatchFileBetweenTags{\JarlskogVal}{constants.tex}{JarlskogVal}
\CatchFileBetweenTags{\JarlskogExperimentalVal}{constants.tex}{JarlskogExperimentalVal}
\CatchFileBetweenTags{\JarlskogSigma}{constants.tex}{JarlskogSigma}

\CatchFileBetweenTags{\VonKlitzingVal}{constants.tex}{VonKlitzingVal}
\CatchFileBetweenTags{\VonKlitzingExperimentalVal}{constants.tex}{VonKlitzingRkExperimentalVal}
\CatchFileBetweenTags{\VonKlitzingSigma}{constants.tex}{VonKlitzingRkSigma}

\CatchFileBetweenTags{\bosonictotalchiVal}{constants.tex}{bosonictotalchiVal}
\CatchFileBetweenTags{\bosonicreducedchiVal}{constants.tex}{bosonicreducedchiVal}
\CatchFileBetweenTags{\spectrumtotalchiVal}{constants.tex}{spectrumtotalchiVal}
\CatchFileBetweenTags{\spectrumreducedchiVal}{constants.tex}{spectrumreducedchiVal}
\CatchFileBetweenTags{\combinedtotalchiVal}{constants.tex}{combinedtotalchiVal}
\CatchFileBetweenTags{\combinedreducedchiVal}{constants.tex}{combinedreducedchiVal}

\CatchFileBetweenTags{\WBosonMassExperimentalVal}{constants.tex}{WBosonMassExperimentalVal}
\CatchFileBetweenTags{\WBosonMassSigma}{constants.tex}{WBosonMassSigma}

\CatchFileBetweenTags{\WeakAngleExperimentalVal}{constants.tex}{WeakAngleExperimentalVal}
\CatchFileBetweenTags{\WeakAngleSigma}{constants.tex}{WeakAngleSigma}

\CatchFileBetweenTags{\ElectronYukawaVal}{constants.tex}{ElectronYukawaVal}
\CatchFileBetweenTags{\ElectronYukawaExperimentalVal}{constants.tex}{ElectronYukawaExperimentalVal}
\CatchFileBetweenTags{\ElectronYukawaSigma}{constants.tex}{ElectronYukawaSigma}

\CatchFileBetweenTags{\GravCouplingVal}{constants.tex}{GravCouplingVal}
\CatchFileBetweenTags{\GravCouplingExperimentalVal}{constants.tex}{GravCouplingExperimentalVal}

\CatchFileBetweenTags{\VacuumEnergyScaleVal}{constants.tex}{VacuumEnergyScaleVal}
\CatchFileBetweenTags{\VacuumEnergyScaleExperimentalVal}{constants.tex}{VacuumEnergyScaleExperimentalVal}

\CatchFileBetweenTags{\HiggsLambdaSigma}{constants.tex}{HiggsLambdaSigma}
\CatchFileBetweenTags{\PlanckMassSigma}{constants.tex}{PlanckMassSigma}
\CatchFileBetweenTags{\GravCouplingSigma}{constants.tex}{GravCouplingSigma}

\CatchFileBetweenTags{\HiggsVEVExperimentalVal}{constants.tex}{HiggsVEVExperimentalVal}
\CatchFileBetweenTags{\HiggsVEVSigma}{constants.tex}{HiggsVEVSigma}

The joint tests of \cref{sec:Synthesis} decide which premises can fail. This section records how the proposed values compare with measurement. A goodness-of-fit summary does not replace those joint tests.

\subsection{A Random-Walk Estimate of the Overhead}
A Monte Carlo comparison of bulk (4D) and surface (3D) random walks on the projected $E_8$ lattice gives a path-length ratio $\eta_{sim}\approx0.988$, within $0.7\%$ of $171/172\approx0.994$. The agreement is suggestive. It is not a derivation of the overhead: the register, reference and event premises identified in \cref{subsec:coordinate_overhead,subsec:weak_mixing} remain open, and the simulation's observable has not been mapped to them. The source code is available as Supplementary Material.

\subsection{Comparison Table}
No coefficient below is fitted to data. Each value is nonetheless conditional on the workload, response and matching premises stated in its section, so the table is not a zero-parameter fit of a completed theory. The W mass uses the measured Z mass and tests the ratio $1-s$; the absolute W/Z response needs the interface factor of \cref{subsec:weak_mixing}. The strong coupling at $M_Z$ uses the zero-momentum impedance. The von Klitzing constant is fixed by $\alpha^{-1}$ through $R_K=\mu_0c/(2\alpha)$ and is not independent of it. Tier~3 values are algebraic consequences of Tier~1 and~2 inputs and are excluded from $\chi^2$.

\begin{table*}[h!]
\centering
\caption{Proposed values compared with measurement. The reduced $\chi^2$ summarizes agreement of separately evaluated formulas; it does not test their joint consistency (\cref{sec:Synthesis}). The electron Yukawa is fixed by the mass law from $v$ and is listed for reference only.}
\renewcommand{\arraystretch}{1.2}
\begin{tabular}{@{}lllr@{}}
\toprule
\textbf{Observable} & \textbf{Proposed Value} & \textbf{Experimental Value} & \textbf{Deviation} \\
\midrule
\multicolumn{4}{c}{\textit{Tier 1: Bosonic Structure \& Couplings}} \\
\midrule
$\alpha^{-1}$ & \AlphaInvVal & \AlphaInvExperimentalVal & $\AlphaInvSigma \sigma$ \\
$G_F$ & \FermiConstVal & \FermiConstExperimentalVal & $\FermiConstSigma \sigma$ \\
$M_W$ ($M_Z$ input) & \WBosonMassVal & \WBosonMassExperimentalVal & $\WBosonMassSigma \sigma$ \\
$\alpha_s(M_Z)$ & \AlphaSVal & \AlphaSExperimentalVal & $\AlphaSSigma \sigma$ \\
$M_H$ & \HiggsMassVal & \HiggsMassExperimentalVal & $\HiggsMassSigma \sigma$ \\
\midrule
\multicolumn{3}{l}{\textbf{Tier 1 $\chi^2$}} & \textbf{$\bosonictotalchiVal$ (5 obs.)} \\
\midrule
\multicolumn{4}{c}{\textit{Tier 2: Mass Hierarchy \& Resistance}} \\
\midrule
$\sin^2\theta_W$ & \WeakAngleVal & \WeakAngleExperimentalVal & $\WeakAngleSigma \sigma$ \\
$J_{CKM}$\textsuperscript{a} & $\JarlskogVal$ & \JarlskogExperimentalVal & $\JarlskogSigma \sigma$ \\
$M_P$ & \PlanckMassVal & \PlanckMassExperimentalVal & $\PlanckMassSigma \sigma$ \\
$R_K$ & \VonKlitzingVal & \VonKlitzingExperimentalVal & $\VonKlitzingSigma \sigma$ \\
\midrule
\multicolumn{3}{l}{\textbf{Tier 2 $\chi^2$}} & \textbf{$\spectrumtotalchiVal$ (4 obs.)} \\
\midrule
\multicolumn{3}{l}{\textbf{Combined Total $\chi^2$}} & \textbf{\combinedtotalchiVal (9 obs.)} \\
\multicolumn{3}{l}{\textbf{Combined Reduced $\chi^2$}} & \textbf{\combinedreducedchiVal} \\
\midrule
\multicolumn{4}{c}{\textit{Tier 3: Derived Consequences}} \\
\midrule
$\lambda$ & $\HiggsLambdaVal$ & \HiggsLambdaExperimentalVal & $\HiggsLambdaSigma \sigma$ \\
$\alpha_G$ & $\GravCouplingVal$ & \GravCouplingExperimentalVal & $\GravCouplingSigma \sigma$ \\
$\rho_{vac}/M_P^4$ & $\VacuumEnergyScaleVal$ & \VacuumEnergyScaleExperimentalVal & Consistent \\
$v$ & $\HiggsVEVVal$ & \HiggsVEVExperimentalVal & $\HiggsVEVSigma \sigma$ \\
$y_e$ (mass law) & $\ElectronYukawaVal$ & \ElectronYukawaExperimentalVal & --- \\
\bottomrule
\label{tab:global_fit}
\end{tabular}

\raggedright\footnotesize\textsuperscript{a}From Paper~IV's derived mixing matrix. The earlier geometric estimate $\phi^2\eta T_{geo}\approx3.085\times10^{-5}$ differs by about $2\%$ and is not used (\cref{subsec:jarlskog}).
\end{table*}

\subsection{Reading the Comparison}
Most proposed values lie within their measurement uncertainties, and the gravitational coupling agrees with $Gm_e^2/\hbar c$ to a few parts per million. These agreements motivate the derivation program; they do not establish it. The decisive tests are the joint relations: the electron mass law, the absolute W and Z masses with one interface response, the strong-coupling running with consistent flavour thresholds, and the scalar interface budget. The W mass in particular separates this proposal ($80.385$~GeV) from standard radiative corrections on the same inputs ($80.356$~GeV) at a level future measurements can resolve.
